\documentclass[12pt, oneside, letterpaper]{book}
\usepackage{configuraciones/setup}

% Casos de USO
% ==============================================================================
% LIBRERÍA DE DIAGRAMAS TIKZ - VERSIÓN CORREGIDA (SIN ENCIMAMIENTOS)
% ==============================================================================

\tikzset{
    actor/.style={thick},
    sistema/.style={thick, fill=gray!5},
    caso/.style={draw, ellipse, fill=azulESCOM!15, minimum width=4.5cm, minimum height=0.8cm, align=center},
    relación/.style={thick},
    extension/.style={thick, dashed, <-, >=stealth}
}

\newcommand{\dibujarQuejoso}[2]{
    \draw[actor] (#1,#2+0.7) circle (0.2); 
    \draw[actor] (#1,#2+0.5) -- (#1,#2); 
    \draw[actor] (#1-0.3,0.3+#2) -- (#1+0.3,0.3+#2); 
    \draw[actor] (#1,#2) -- (#1-0.2,#2-0.3); 
    \draw[actor] (#1,#2) -- (#1+0.2,#2-0.3); 
    \node at (#1,#2-0.6) {\scriptsize Quejoso};
}

% --- FIGURAS CORREGIDAS ---

% MODULO PERFIL
\newcommand{\figCUuno}{ \begin{tikzpicture}[font=\small] \dibujarQuejoso{-4}{2} \draw[sistema] (-1.5,0.5) rectangle (6,3.5); \node[anchor=north] at (2.25,3.4) {\textbf{\scriptsize Gestión de Perfil}}; \node[caso] (c) at (2.25,2) {\scriptsize CU01: Crear cuenta de usuario}; \draw[relación] (-3.7,2.3) -- (c.west); \end{tikzpicture} }

\newcommand{\figCUdos}{ \begin{tikzpicture}[font=\small] \dibujarQuejoso{-4}{2} \draw[sistema] (-1.5,0.5) rectangle (6,3.5); \node[anchor=north] at (2.25,3.4) {\textbf{\scriptsize Gestión de Perfil}}; \node[caso] (c) at (2.25,2) {\scriptsize CU02: Iniciar sesión}; \draw[relación] (-3.7,2.3) -- (c.west); \end{tikzpicture} }

\newcommand{\figCUtres}{ \begin{tikzpicture}[font=\small] \dibujarQuejoso{-4}{2.5} \draw[sistema] (-1.5,-0.5) rectangle (6,5); \node[caso] (c2) at (2.25,4) {\scriptsize CU02: Iniciar sesión}; \node[caso] (c3) at (2.25,0.5) {\scriptsize CU03: Recuperar contraseña}; \draw[extension] (c2) -- (c3) node[midway, right=0.2cm] {\tiny $<<Extend>>$}; \draw[relación] (-3.7,2.8) -- (c2.west); \end{tikzpicture} }

\newcommand{\figCUcuatro}{ \begin{tikzpicture}[font=\small] \dibujarQuejoso{-4}{2} \draw[sistema] (-1.5,0.5) rectangle (6,3.5); \node[caso] (c) at (2.25,2) {\scriptsize CU04: Editar perfil de usuario}; \draw[relación] (-3.7,2.3) -- (c.west); \end{tikzpicture} }

\newcommand{\figCUcinco}{ \begin{tikzpicture}[font=\small] \dibujarQuejoso{-4}{2} \draw[sistema] (-1.5,0.5) rectangle (6,3.5); \node[caso] (c) at (2.25,2) {\scriptsize CU05: Registrar datos de tutor}; \draw[relación] (-3.7,2.3) -- (c.west); \end{tikzpicture} }

% MODULO QUEJAS
\newcommand{\figCUseis}{ \begin{tikzpicture}[font=\small] \dibujarQuejoso{-4}{2} \draw[sistema] (-1.5,0.5) rectangle (6,3.5); \node[caso] (c) at (2.25,2) {\scriptsize CU06: Levantar queja}; \draw[relación] (-3.7,2.3) -- (c.west); \end{tikzpicture} }

\newcommand{\figCUsiete}{ \begin{tikzpicture}[font=\small] \dibujarQuejoso{-4}{2} \draw[sistema] (-1.5,0.5) rectangle (6,3.5); \node[caso] (c) at (2.25,2) {\scriptsize CU07: Adjuntar evidencia}; \draw[relación] (-3.7,2.3) -- (c.west); \end{tikzpicture} }

\newcommand{\figCUocho}{ \begin{tikzpicture}[font=\small] \dibujarQuejoso{-4}{2} \draw[sistema] (-1.5,0.5) rectangle (6,3.5); \node[caso] (c) at (2.25,2) {\scriptsize CU08: Consultar tablero de trámites}; \draw[relación] (-3.7,2.3) -- (c.west); \end{tikzpicture} }

% FIGURAS CON EXTENDS (CORREGIDAS)
\newcommand{\figCUnueve}{ \begin{tikzpicture}[font=\small] \dibujarQuejoso{-4}{2} \node[caso] (c6) at (0,2) {\scriptsize CU06: Levantar queja}; \node[caso] (c9) at (6,2) {\scriptsize CU09: Modificar queja}; \draw[extension] (c6) -- (c9) node[midway, above=0.1cm] {\tiny $<<Extend>>$}; \draw[relación] (-3.7,2.3) -- (c6.west); \end{tikzpicture} }

\newcommand{\figCUdiez}{ \begin{tikzpicture}[font=\small] \dibujarQuejoso{-4}{2} \node[caso] (c6) at (0,2) {\scriptsize CU06: Levantar queja}; \node[caso] (c10) at (6,2) {\scriptsize CU10: Eliminar borrador}; \draw[extension] (c6) -- (c10) node[midway, above=0.1cm] {\tiny $<<Extend>>$}; \draw[relación] (-3.7,2.3) -- (c6.west); \end{tikzpicture} }

\newcommand{\figCUonce}{ \begin{tikzpicture}[font=\small] \dibujarQuejoso{-4}{2} \draw[sistema] (-1.5,0.5) rectangle (6,3.5); \node[caso] (c) at (2.25,2) {\scriptsize CU11: Recibir notificaciones}; \draw[relación] (-3.7,2.3) -- (c.west); \end{tikzpicture} }

\newcommand{\figCUdoce}{ \begin{tikzpicture}[font=\small] \dibujarQuejoso{-4}{2} \node[caso] (c11) at (0,2) {\scriptsize CU11: Notificaciones}; \node[caso] (c12) at (6,2) {\scriptsize CU12: Completar información}; \draw[extension] (c11) -- (c12) node[midway, above=0.1cm] {\tiny $<<Extend>>$}; \draw[relación] (-3.7,2.3) -- (c11.west); \end{tikzpicture} }

\newcommand{\figCUtrece}{ \begin{tikzpicture}[font=\small] \dibujarQuejoso{-4}{2} \node[caso] (c11) at (0,2) {\scriptsize CU11: Notificaciones}; \node[caso] (c13) at (6,2) {\scriptsize CU13: Medidas precautorias}; \draw[extension] (c11) -- (c13) node[midway, above=0.1cm] {\tiny $<<Extend>>$}; \draw[relación] (-3.7,2.3) -- (c11.west); \end{tikzpicture} }

\newcommand{\figCUcatorce}{ \begin{tikzpicture}[font=\small] \dibujarQuejoso{-4}{2} \node[caso] (c11) at (0,2) {\scriptsize CU11: Notificaciones}; \node[caso] (c14) at (6,2) {\scriptsize CU14: Propuesta conciliación}; \draw[extension] (c11) -- (c14) node[midway, above=0.1cm] {\tiny $<<Extend>>$}; \draw[relación] (-3.7,2.3) -- (c11.west); \end{tikzpicture} }

\newcommand{\figCUquince}{ \begin{tikzpicture}[font=\small] \dibujarQuejoso{-4}{2} \node[caso] (c14) at (0,2) {\scriptsize CU14: Propuesta}; \node[caso] (c15) at (6,2) {\scriptsize CU15: Aceptar propuesta}; \draw[extension] (c14) -- (c15) node[midway, above=0.1cm] {\tiny $<<Extend>>$}; \draw[relación] (-3.7,2.3) -- (c14.west); \end{tikzpicture} }

\newcommand{\figCUdieciseis}{ \begin{tikzpicture}[font=\small] \dibujarQuejoso{-4}{2} \node[caso] (c14) at (0,2) {\scriptsize CU14: Propuesta}; \node[caso] (c16) at (6,2) {\scriptsize CU16: Rechazar propuesta}; \draw[extension] (c14) -- (c16) node[midway, above=0.1cm] {\tiny $<<Extend>>$}; \draw[relación] (-3.7,2.3) -- (c14.west); \end{tikzpicture} }

\newcommand{\figCUdiecisiete}{ \begin{tikzpicture}[font=\small] \dibujarQuejoso{-4}{2} \draw[sistema] (-1.5,0.5) rectangle (6,3.5); \node[caso] (c) at (2.25,2) {\scriptsize CU17: Descargar acuse}; \draw[relación] (-3.7,2.3) -- (c.west); \end{tikzpicture} }

\newcommand{\figCUdieciocho}{ \begin{tikzpicture}[font=\small] \dibujarQuejoso{-4}{2} \draw[sistema] (-1.5,0.5) rectangle (6,3.5); \node[caso] (c) at (2.25,2) {\scriptsize CU18: Detalles completos}; \draw[relación] (-3.7,2.3) -- (c.west); \end{tikzpicture} }
\usepackage[utf8]{inputenc}
\usepackage{textcomp}
\lstset{inputencoding=utf8, extendedchars=true}

% Imágenes y PDFs
\usepackage{pdfpages} % Permite insertar páginas de archivos PDF externos
\usepackage{ragged2e}

% Definición de estilo para JSON
\lstdefinelanguage{json}{
    basicstyle=\normalfont\ttfamily\small,
    numbers=left,
    numberstyle=\scriptsize,
    stepnumber=1,
    numbersep=8pt,
    showstringspaces=false,
    breaklines=true,
    frame=lines,
    backgroundcolor=\color{gray!5}
}

\usepackage{tabularx} % Para que las tablas ocupen el ancho de la página
\usepackage{enumitem} % Para reducir espacio en listas si las usas
\usepackage[table]{xcolor} % Por si quieres dar color a las cabeceras

% ==============================================================================
\usepackage{hyperref} 
\hypersetup{
    colorlinks=true,
    linkcolor=blue,
    urlcolor=blue,
    pdftitle={Trabajo Terminal - Defensoría de los Derechos Politécnicos},
    pdfnewwindow=true % Esto ayuda a que los archivos externos se abran en otra ventana
}

\begin{document}

\section{Catálogo de Requerimientos y Reglas de Negocio}

Este documento presenta el catálogo completo de requerimientos y reglas de negocio
levantados durante la etapa de interacción con el cliente. Los artefactos se
organizan en cinco categorías: requerimientos de usuario (RU), requerimientos de
sistema (RS), requerimientos funcionales (RF), requerimientos no funcionales (RNF)
y reglas de negocio (RN).

%============================================================
\subsection{Requerimientos de Usuario}
%============================================================

En esta sección se detallan las necesidades específicas de cada actor que
interactúa con la plataforma. Estos requerimientos representan las acciones
fundamentales que el sistema debe permitir para garantizar el ciclo de vida
de la queja y el cumplimiento normativo institucional.

\begin{table}[H]
\centering
\caption{Requerimientos de Usuario Específicos por Actor.}
\begin{tabularx}{\textwidth}{@{} l l L @{}}
\toprule
\rowcolor[HTML]{2D3E50}
\textbf{\color{white} ID} & \textbf{\color{white} Actor} &
\textbf{\color{white} Descripción de la Necesidad} \\ \midrule
\rowcolor[HTML]{F2F4F7}
RU-01 & Quejoso & Registro de cuenta mediante credenciales institucionales (boleta/empleado). \\
RU-02 & Quejoso & Gestión de perfil personal y carga de datos de tutor (en caso de ser menor de edad). \\
\rowcolor[HTML]{F2F4F7}
RU-03 & Quejoso & Levantamiento de queja detallando hechos, lugar y autoridades involucradas. \\
RU-04 & Quejoso & Carga de evidencias multimedia (PDF, imágenes, audio, video) para sustentar la queja. \\
\rowcolor[HTML]{F2F4F7}
RU-05 & Quejoso & Consulta del estatus de la queja en tiempo real mediante un folio digital único. \\
RU-06 & Quejoso & Recepción de notificaciones automáticas sobre actualizaciones en su expediente. \\ \midrule
\rowcolor[HTML]{F2F4F7}
RU-07 & Recepcionista & Visualización de quejas entrantes en una bandeja de entrada centralizada. \\
RU-08 & Recepcionista & Validación de requisitos mínimos y triaje inicial de la documentación recibida. \\ \midrule
\rowcolor[HTML]{F2F4F7}
RU-09 & Analista & Evaluación jurídica de fondo de los hechos narrados para determinar competencia. \\
RU-10 & Analista & Gestión del expediente digital, incluyendo la adición de notas y documentos internos. \\ \midrule
\rowcolor[HTML]{F2F4F7}
RU-11 & Subdefensor & Supervisión de los plazos legales para evitar el rezago administrativo. \\
RU-12 & Subdefensor & Gestión de actos de investigación y seguimiento de oficios enviados a autoridades. \\ \midrule
\rowcolor[HTML]{F2F4F7}
RU-13 & Defensor & Visualización de estadísticas generales y tableros de control de gestión institucional. \\
RU-14 & Defensor & Emisión, firma digital avanzada y validación de acuerdos de conclusión del caso. \\
\rowcolor[HTML]{F2F4F7}
RU-15 & Defensor & Administración de privilegios de usuario y configuración de roles del sistema (RBAC). \\ \bottomrule
\end{tabularx}
\end{table}


%============================================================
\subsection{Requerimientos de Sistema}
%============================================================

En esta sección se definen las capacidades técnicas y funcionales que el sistema
debe ejecutar para dar soporte a los procesos operativos de la Defensoría. Estos
requerimientos están enfocados en la lógica de procesamiento, seguridad de datos
y automatización de trámites.

\begin{table}[H]
\centering
\caption{Requerimientos de Sistema (Funcionales y Técnicos).}
\begin{tabularx}{\textwidth}{@{} l L l @{}}
\toprule
\rowcolor[HTML]{2D3E50}
\textbf{\color{white} Categoría} & \textbf{\color{white} Especificación Técnica} &
\textbf{\color{white} Alcance} \\ \midrule
\rowcolor[HTML]{F2F4F7}
Inteligencia Artificial & Motor de clasificación basado en PLN para determinar la competencia de las quejas de forma automática. & Sistema \\
Seguridad (RBAC) & Sistema de autenticación con JWT y control de acceso basado en roles para restringir la visibilidad de datos sensibles. & General \\
\rowcolor[HTML]{F2F4F7}
Gestión Documental & Generador automático de plantillas para oficios a Unidades Académicas, acuerdos de admisión y de conclusión en PDF. & Subdefensor \\
Firma Electrónica & Módulo de sellado digital y firma avanzada para garantizar la integridad y validez jurídica de los acuerdos finales. & Defensor \\
\rowcolor[HTML]{F2F4F7}
Gestión de Plazos & Lógica de monitoreo de términos procesales (2 a 15 días) con alertas automáticas por vencimiento de plazos legales. & Subdefensor \\
Notificaciones & Servicio de mensajería automática (email/notificación push) para informar al quejoso cambios en el estatus de su folio. & General \\
\rowcolor[HTML]{F2F4F7}
Búsqueda Avanzada & Motor de indexación GIN para realizar búsquedas de texto completo (\textit{full-text search}) dentro del repositorio. & Analista \\
Trazabilidad & Sistema de registro de auditoría (\textit{logs}) para rastrear cada modificación, acceso o cambio de estado en el expediente. & Administrador \\
\rowcolor[HTML]{F2F4F7}
Analítica de Datos & Dashboard interactivo con generación de métricas sobre reincidencia de autoridades y tipos de derechos vulnerados. & Defensor \\
Cifrado en Reposo & Implementación de \texttt{pgcrypto} (AES-256) para el almacenamiento cifrado de datos personales e imágenes de evidencia. & Sistema \\ \bottomrule
\end{tabularx}
\end{table}


%============================================================
\subsection{Requerimientos Funcionales}
%============================================================

A partir del análisis de los 75 casos de uso definidos para la plataforma,
se extraen los siguientes requerimientos funcionales desglosados por módulos
operativos.

\begin{table}[H]
\centering
\caption{RF — Módulo de Acceso y Gestión de Identidad.}
\begin{tabularx}{\textwidth}{@{} l L l @{}}
\toprule
\rowcolor[HTML]{2D3E50}
\textbf{\color{white} ID} & \textbf{\color{white} Descripción} & \textbf{\color{white} Actor} \\ \midrule
\rowcolor[HTML]{F2F4F7}
RF-01 & Registro de usuarios con validación de identidad institucional (Boleta/Número de Empleado). & Quejoso \\
RF-02 & Recuperación de contraseña mediante el envío de tokens temporales al correo registrado. & Usuario \\
\rowcolor[HTML]{F2F4F7}
RF-03 & Gestión de perfil de usuario permitiendo la actualización de datos de contacto y domicilio. & Quejoso \\
RF-04 & Captura obligatoria de datos del tutor y documentos de identidad para quejosos menores de edad. & Quejoso \\
\rowcolor[HTML]{F2F4F7}
RF-05 & Autenticación de personal administrativo mediante roles definidos en el sistema (RBAC). & Personal DDP \\ \bottomrule
\end{tabularx}
\end{table}

\begin{table}[H]
\centering
\caption{RF — Módulo de Gestión de Quejas (\textit{Front-end}).}
\begin{tabularx}{\textwidth}{@{} l L l @{}}
\toprule
\rowcolor[HTML]{2D3E50}
\textbf{\color{white} ID} & \textbf{\color{white} Descripción} & \textbf{\color{white} Actor} \\ \midrule
\rowcolor[HTML]{F2F4F7}
RF-06 & Formulario dinámico para el levantamiento de quejas con narrativa de hechos y geolocalización. & Quejoso \\
RF-07 & Carga de evidencias en múltiples formatos (PDF, JPG, MP3, MP4) con límite de tamaño por archivo. & Quejoso \\
\rowcolor[HTML]{F2F4F7}
RF-08 & Generación de folio único con nomenclatura institucional tras finalizar el registro exitoso. & Sistema \\
RF-09 & Visualización de tablero personal con el listado de quejas activas y su estatus actual. & Quejoso \\
\rowcolor[HTML]{F2F4F7}
RF-10 & Consultar línea de tiempo interactiva que detalle el historial de pasos de la queja. & Quejoso \\
RF-11 & Descarga de acuse de recibo y documentos oficiales en formato PDF. & Quejoso \\ \bottomrule
\end{tabularx}
\end{table}

\begin{table}[H]
\centering
\caption{RF — Módulo de Recepción y Administrativo.}
\begin{tabularx}{\textwidth}{@{} l L l @{}}
\toprule
\rowcolor[HTML]{2D3E50}
\textbf{\color{white} ID} & \textbf{\color{white} Descripción} & \textbf{\color{white} Actor} \\ \midrule
\rowcolor[HTML]{F2F4F7}
RF-12 & Bandeja de entrada centralizada para la recepción de solicitudes con filtros de prioridad. & Recepcionista \\
RF-13 & Búsqueda de antecedentes históricos para detectar duplicidad o reincidencia de quejas. & Recepcionista \\
\rowcolor[HTML]{F2F4F7}
RF-14 & Emisión de notificaciones de prevención para solicitar información faltante al quejoso. & Recepcionista \\
RF-15 & Función para rechazar quejas no competentes con justificación legal obligatoria. & Recepcionista \\
\rowcolor[HTML]{F2F4F7}
RF-16 & Turnar expedientes validados al área de análisis para el inicio del proceso formal. & Recepcionista \\ \bottomrule
\end{tabularx}
\end{table}

\begin{table}[H]
\centering
\caption{RF — Módulo de Análisis.}
\begin{tabularx}{\textwidth}{@{} l L l @{}}
\toprule
\rowcolor[HTML]{2D3E50}
\textbf{\color{white} ID} & \textbf{\color{white} Descripción} & \textbf{\color{white} Actor} \\ \midrule
\rowcolor[HTML]{F2F4F7}
RF-17 & Determinación automática de competencia preliminar mediante el motor de IA. & Sistema \\
RF-18 & Adición de notas internas y comentarios técnicos visibles solo para personal jurídico. & Analista \\
\rowcolor[HTML]{F2F4F7}
RF-19 & Generación de Acuerdo de Admisión vinculando los hechos con la normativa del IPN. & Analista \\
RF-20 & Asignación de expediente a un abogado específico para la etapa de investigación. & Subdefensor \\
\rowcolor[HTML]{F2F4F7}
RF-21 & Registro y control de solicitudes de actos de investigación dirigidos a autoridades de Unidades Académicas. & Subdefensor \\
RF-22 & Gestión de plazos con semáforo de alertas (Verde, Amarillo, Rojo) según el vencimiento. & Subdefensor \\
\rowcolor[HTML]{F2F4F7}
RF-23 & Redacción y envío de recordatorios automáticos a autoridades que no han rendido informes. & Subdefensor \\ \bottomrule
\end{tabularx}
\end{table}

\begin{table}[H]
\centering
\caption{RF — Módulo de Resolución, Firma y Cierre.}
\begin{tabularx}{\textwidth}{@{} l L l @{}}
\toprule
\rowcolor[HTML]{2D3E50}
\textbf{\color{white} ID} & \textbf{\color{white} Descripción} & \textbf{\color{white} Actor} \\ \midrule
\rowcolor[HTML]{F2F4F7}
RF-24 & Elaboración de propuestas de conciliación y envío para aceptación o rechazo del quejoso. & Analista \\
RF-25 & Consultar tableros de gestión global con estadísticas de eficiencia operativa de la DDP. & Defensor \\
\rowcolor[HTML]{F2F4F7}
RF-26 & Ejecución de Firma Digital Avanzada sobre el Acuerdo de Conclusión individual o masiva. & Defensor \\
RF-27 & Actualización automática del estado del expediente a ``FINALIZADO'' tras la firma digital. & Sistema \\
\rowcolor[HTML]{F2F4F7}
RF-28 & Transferencia de expedientes concluidos al repositorio de archivo histórico digital. & Sistema \\ \bottomrule
\end{tabularx}
\end{table}

\begin{table}[H]
\centering
\caption{RF — Módulo de Administración Técnica.}
\begin{tabularx}{\textwidth}{@{} l L l @{}}
\toprule
\rowcolor[HTML]{2D3E50}
\textbf{\color{white} ID} & \textbf{\color{white} Descripción} & \textbf{\color{white} Actor} \\ \midrule
\rowcolor[HTML]{F2F4F7}
RF-29 & Configuración de plantillas dinámicas para oficios, acuerdos y notificaciones institucionales. & Administrador \\
RF-30 & Gestión del catálogo de dependencias politécnicas y autoridades registradas en el sistema. & Administrador \\ \bottomrule
\end{tabularx}
\end{table}


%============================================================
\subsection{Requerimientos No Funcionales}
%============================================================

Los requerimientos no funcionales definen los atributos de calidad,
restricciones técnicas y estándares de seguridad que el sistema debe cumplir
para garantizar una operación robusta, confiable y alineada con la normativa
institucional.

% Tabla 1 — Seguridad y Privacidad
% Originales conservados: RNF-01, RNF-02, RNF-05
% Eliminados: RNF-03 (AES-256), RNF-04 (TLS), RNF-06 (no repudio)
% Renumerados: RNF-05 → RNF-03
\begin{table}[H]
\centering
\caption{RNF — Atributos de Seguridad y Privacidad.}
\begin{tabularx}{\textwidth}{@{} l L l @{}}
\toprule
\rowcolor[HTML]{2D3E50}
\textbf{\color{white} ID} & \textbf{\color{white} Descripción del Atributo} & \textbf{\color{white} Categoría} \\ \midrule
\rowcolor[HTML]{F2F4F7}
RNF-01 & El acceso administrativo debe estar restringido mediante el modelo RBAC (Role-Based Access Control). & Seguridad \\
RNF-02 & Implementación de tokens JWT con tiempo de expiración para la gestión de sesiones seguras. & Seguridad \\
\rowcolor[HTML]{F2F4F7}
RNF-03 & Anonimización de datos sensibles en el motor de IA para proteger la identidad durante el entrenamiento. & Privacidad \\ \bottomrule
\end{tabularx}
\end{table}

% Tabla 2 — Rendimiento y Escalabilidad
% Originales conservados: RNF-07, RNF-09, RNF-10
% Eliminados: RNF-08, RNF-11
% Renumerados: RNF-07 → RNF-04 | RNF-09 → RNF-05 | RNF-10 → RNF-06
\begin{table}[H]
\centering
\caption{RNF — Atributos de Rendimiento y Escalabilidad.}
\begin{tabularx}{\textwidth}{@{} l L l @{}}
\toprule
\rowcolor[HTML]{2D3E50}
\textbf{\color{white} ID} & \textbf{\color{white} Descripción del Atributo} & \textbf{\color{white} Categoría} \\ \midrule
\rowcolor[HTML]{F2F4F7}
RNF-04 & El tiempo de respuesta para consultas de estatus no debe superar los 2 segundos bajo carga normal. & Rendimiento \\
RNF-05 & La clasificación automática de quejas por el motor de IA debe realizarse en un tiempo menor a 5 segundos. & Rendimiento \\
\rowcolor[HTML]{F2F4F7}
RNF-06 & Capacidad de escalado horizontal de microservicios mediante orquestación con Docker Compose o Kubernetes. & Escalabilidad \\ \bottomrule
\end{tabularx}
\end{table}

% Tabla 3 — Fiabilidad y Disponibilidad
% Originales conservados: RNF-13, RNF-14, RNF-15
% Eliminados: RNF-12
% Renumerados: RNF-13 → RNF-07 | RNF-14 → RNF-08 | RNF-15 → RNF-09
\begin{table}[H]
\centering
\caption{RNF — Atributos de Fiabilidad y Disponibilidad.}
\begin{tabularx}{\textwidth}{@{} l L l @{}}
\toprule
\rowcolor[HTML]{2D3E50}
\textbf{\color{white} ID} & \textbf{\color{white} Descripción del Atributo} & \textbf{\color{white} Categoría} \\ \midrule
\rowcolor[HTML]{F2F4F7}
RNF-07 & Tolerancia a fallos: la caída del microservicio de IA no debe afectar la gestión administrativa de quejas. & Resiliencia \\
RNF-08 & Integridad de los datos: las transacciones en la base de datos deben cumplir estrictamente con el modelo ACID. & Fiabilidad \\
\rowcolor[HTML]{F2F4F7}
RNF-09 & El sistema debe realizar respaldos automáticos diarios (\textit{backups}) del repositorio digital de evidencias. & Recuperación \\ \bottomrule
\end{tabularx}
\end{table}

% Tabla 4 — Mantenibilidad y Portabilidad
% Originales conservados: RNF-16, RNF-17, RNF-18
% Eliminados: RNF-19
% Renumerados: RNF-16 → RNF-10 | RNF-17 → RNF-11 | RNF-18 → RNF-12
\begin{table}[H]
\centering
\caption{RNF — Atributos de Mantenibilidad y Portabilidad.}
\begin{tabularx}{\textwidth}{@{} l L l @{}}
\toprule
\rowcolor[HTML]{2D3E50}
\textbf{\color{white} ID} & \textbf{\color{white} Descripción del Atributo} & \textbf{\color{white} Categoría} \\ \midrule
\rowcolor[HTML]{F2F4F7}
RNF-10 & Arquitectura desacoplada de microservicios para permitir actualizaciones modulares sin afectar todo el sistema. & Mantenibilidad \\
RNF-11 & Código fuente documentado y gestionado mediante control de versiones en GitHub para facilitar la colaboración. & Mantenibilidad \\
\rowcolor[HTML]{F2F4F7}
RNF-12 & Contenerización de todos los servicios mediante Docker para asegurar la portabilidad entre entornos. & Portabilidad \\ \bottomrule
\end{tabularx}
\end{table}


%============================================================
\subsection{Reglas de Negocio}
%============================================================

Las reglas de negocio establecen las restricciones y directrices operativas
que el sistema debe validar para garantizar que los procesos administrativos
se realicen bajo los principios de legalidad, transparencia y seguridad
institucional de la Defensoría.

\begin{table}[H]
\centering
\caption{RN — Acceso, Seguridad y Auditoría.}
\begin{tabularx}{\textwidth}{@{} l X @{}}
\toprule
\rowcolor[HTML]{2D3E50}
\textbf{\color{white} ID} & \textbf{\color{white} Regla de Negocio} \\ \midrule
\rowcolor[HTML]{F2F4F7}
RN-01 & Solo miembros activos de la comunidad politécnica con boleta o número de empleado pueden registrar quejas. \\
RN-02 & El acceso administrativo es restringido por roles; un actor no puede ejecutar funciones de un nivel superior (RBAC). \\
\rowcolor[HTML]{F2F4F7}
RN-03 & La asignación, alta o modificación de permisos de usuario debe registrarse en un \textit{log} de auditoría inmutable. \\
RN-04 & Los datos personales del quejoso deben permanecer cifrados en la base de datos y solo ser visibles por personal autorizado. \\
\rowcolor[HTML]{F2F4F7}
RN-05 & Ningún usuario administrativo puede eliminar registros de auditoría o historiales de transacciones del sistema. \\
RN-06 & El sistema debe cerrar la sesión automáticamente tras 15 minutos de inactividad para prevenir accesos no autorizados. \\ \bottomrule
\end{tabularx}
\end{table}

\begin{table}[H]
\centering
\caption{RN — Registro y Determinación de Competencia.}
\begin{tabularx}{\textwidth}{@{} l X @{}}
\toprule
\rowcolor[HTML]{2D3E50}
\textbf{\color{white} ID} & \textbf{\color{white} Regla de Negocio} \\ \midrule
\rowcolor[HTML]{F2F4F7}
RN-07 & Toda queja debe incluir obligatoriamente una narrativa de hechos, autoridad imputada y unidad académica. \\
RN-08 & Si el quejoso es menor de edad, el sistema exige el registro de un tutor con identificación oficial válida. \\
\rowcolor[HTML]{F2F4F7}
RN-09 & No se aceptarán quejas anónimas; la identidad del promovente debe ser validada contra bases de datos del IPN. \\
RN-10 & Las quejas presentadas después de 90 días naturales de ocurrido el hecho se clasificarán como extemporáneas. \\
\rowcolor[HTML]{F2F4F7}
RN-11 & El sistema no admitirá quejas de carácter laboral o sindical según lo estipulado en el Artículo 13. \\
RN-12 & El motor de IA solo proporciona una sugerencia de competencia; la decisión final es siempre humana. \\
\rowcolor[HTML]{F2F4F7}
RN-13 & La generación del folio único es automática e irreversible; no puede existir duplicidad de folios en el sistema. \\ \bottomrule
\end{tabularx}
\end{table}

\begin{table}[H]
\centering
\caption{RN — Proceso, Sustanciación e Investigación.}
\begin{tabularx}{\textwidth}{@{} l X @{}}
\toprule
\rowcolor[HTML]{2D3E50}
\textbf{\color{white} ID} & \textbf{\color{white} Regla de Negocio} \\ \midrule
\rowcolor[HTML]{F2F4F7}
RN-14 & El estatus ``En Análisis'' se activa únicamente cuando el Recepcionista turna el expediente a la Subdefensoría. \\
RN-15 & El sistema debe bloquear el avance del folio si no se ha cargado el Acuerdo de Admisión correspondiente. \\
\rowcolor[HTML]{F2F4F7}
RN-16 & Todo expediente radicado debe tener asignado obligatoriamente a un abogado analista responsable. \\
RN-17 & Los oficios enviados a Unidades Académicas deben generar automáticamente un recordatorio de respuesta. \\
\rowcolor[HTML]{F2F4F7}
RN-18 & El plazo de respuesta otorgado a la autoridad debe ser de 2 a 15 días hábiles según la normativa. \\
RN-19 & No se puede pasar a la fase de resolución si existen actos de investigación con plazos de respuesta vigentes. \\
\rowcolor[HTML]{F2F4F7}
RN-20 & Las notas internas del Analista no son descargables ni visibles para el quejoso en el portal público. \\
RN-21 & El estatus ``En Investigación'' es obligatorio antes de emitir cualquier propuesta de conclusión de fondo. \\ \bottomrule
\end{tabularx}
\end{table}

\begin{table}[H]
\centering
\caption{RN — Conclusión, Firma y Archivo.}
\begin{tabularx}{\textwidth}{@{} l X @{}}
\toprule
\rowcolor[HTML]{2D3E50}
\textbf{\color{white} ID} & \textbf{\color{white} Regla de Negocio} \\ \midrule
\rowcolor[HTML]{F2F4F7}
RN-22 & El Acuerdo de Conclusión debe ser visado por el Subdefensor antes de enviarse a firma del Titular. \\
RN-23 & La notificación al quejoso de la propuesta de conclusión es requisito previo para el cierre del folio. \\
\rowcolor[HTML]{F2F4F7}
RN-24 & La Firma Digital Avanzada del Defensor es el único medio para validar legalmente el cierre del caso. \\
RN-25 & El sistema impedirá la modificación de cualquier documento una vez que haya sido firmado digitalmente. \\
\rowcolor[HTML]{F2F4F7}
RN-26 & El estatus ``Finalizado'' solo se activa mediante la aplicación exitosa de la firma electrónica del Defensor. \\
RN-27 & Todo expediente finalizado debe ser transferido al archivo histórico en un periodo no mayor a 24 horas. \\
\rowcolor[HTML]{F2F4F7}
RN-28 & El envío a archivo histórico es irreversible y bloquea cualquier edición posterior de los datos del caso. \\ \bottomrule
\end{tabularx}
\end{table}

\begin{table}[H]
\centering
\caption{RN — Estadísticas y Administración.}
\begin{tabularx}{\textwidth}{@{} l X @{}}
\toprule
\rowcolor[HTML]{2D3E50}
\textbf{\color{white} ID} & \textbf{\color{white} Regla de Negocio} \\ \midrule
\rowcolor[HTML]{F2F4F7}
RN-29 & Solo el rol de Defensor (Titular) tiene facultad para visualizar tableros de control global y estadísticas. \\
RN-30 & Las métricas de reincidencia deben calcularse automáticamente basándose en los datos históricos del repositorio. \\
\rowcolor[HTML]{F2F4F7}
RN-31 & La configuración de plantillas de documentos oficiales es de acceso restringido al perfil Administrador. \\
RN-32 & El sistema debe generar alertas preventivas ante el rezago administrativo (expedientes sin movimiento en 10 días). \\
\rowcolor[HTML]{F2F4F7}
RN-33 & Las copias de seguridad de los expedientes deben almacenarse en un repositorio inmutable y externo al sistema. \\
RN-34 & Toda actualización de estado del folio debe generar una notificación automática al correo del quejoso. \\
\rowcolor[HTML]{F2F4F7}
RN-35 & El catálogo de Unidades Académicas y dependencias debe estar sincronizado con la estructura oficial del IPN. \\
RN-36 & No se pueden emitir oficios si la autoridad responsable no se encuentra debidamente registrada en el sistema. \\
\rowcolor[HTML]{F2F4F7}
RN-37 & Los acuerdos de prevención deben otorgar al quejoso un plazo improrrogable de 5 días para subsanar omisiones. \\
RN-38 & El rechazo de una queja por incompetencia debe estar fundado obligatoriamente en los criterios del Artículo 13. \\
\rowcolor[HTML]{F2F4F7}
RN-39 & Las propuestas de conciliación requieren la firma de aceptación del quejoso cargada en el sistema para proceder. \\
RN-40 & El sistema debe garantizar que el folio digital se mantenga accesible para consulta del quejoso por un año mínimo. \\ \bottomrule
\end{tabularx}
\end{table}

\end{document}